% GENERATED FILE. Edit canonical lessons, then run npm run generate:books.
% canonical-source-hash: fnv1a64:0d6efa92ac0baa1d
% canonical-lessons: FA-C509-pazel, FA-C509-takhte-nard, FA-C509-varaq, FA-C509-tas, FA-C509-varaq-bazi-kardan

\chapter{Games}
\label{ch:fa-games}
\hlchaptermodality{509}

\begin{chapteropening}
\textbf{By the end of this chapter:} \emph{I can talk about board and card games.}

Complete the last lesson of chapter 509: I can talk about board and card games.
\end{chapteropening}

\section[\texorpdfstring{\fa{پازل} (pâzel)}{پازل (pâzel)}]{\fa{پازل} (pâzel) — a jigsaw puzzle}

\label{lesson:FA-C509-pazel}



\begin{quote}
Before the new one: say the Persian for a handicraft, then the Persian for pottery.
\end{quote}

\subsection*{You'll want to know: \fa{پازل}}
\textbf{\fa{پازل}} (\emph{pâzel}) — \textquotedblleft{}a jigsaw puzzle\textquotedblright{}.

\begin{quote}
\fa{ما} \fa{شبها} \fa{با} \fa{هم} \fa{پازل} \fa{درست} \fa{میکنیم}.
\end{quote}

\begin{quote}
\emph{mâ shabhâ bâ ham pâzel dorost mikonim.}
\end{quote}

\begin{quote}
We do jigsaws together in the evenings.
\end{quote}

\emph{pâzel dorost kardan} — to do a jigsaw.

\begin{cousinweb}
\emph{pâzel} is a modern loan, usually traced to English \emph{puzzle}. EtymDB has no entry for it.
\end{cousinweb}

\begin{grammarlens}[title={What you've built}]
One more word for telling, explaining and describing at B1.
\end{grammarlens}

\subsection*{Your turn}
\begin{itemize}
\raggedright
  \item \emph{Say it:} \emph{pâzel}
  \item \emph{Say it:} \emph{pâzel}, once more
  \item \emph{Say it:} a sentence of your own with \emph{pâzel}
  \item \emph{Recall:} say the Persian for a handicraft, then the Persian for pottery, then say \emph{pâzel} again
\end{itemize}

\begin{tcolorbox}[breakable,colback=teal!4,colframe=teal!35!black,title={Before you move on}]
What does \fa{پازل} mean? (\textquotedblleft{}A jigsaw puzzle\textquotedblright{}.) Say it once more.
\end{tcolorbox}
\section[\texorpdfstring{\fa{تخته} \fa{نرد} (takhte nard)}{تخته نرد (takhte nard)}]{\fa{تخته} \fa{نرد} (takhte nard) — backgammon}

\label{lesson:FA-C509-takhte-nard}



\begin{quote}
Before the new one: say the Persian for pottery, then the Persian for a jigsaw puzzle.
\end{quote}

\subsection*{You'll want to know: \fa{تخته} \fa{نرد}}
\textbf{\fa{تخته} \fa{نرد}} (\emph{takhte nard}) — \textquotedblleft{}backgammon\textquotedblright{}.

\begin{quote}
\fa{پدربزرگ} \fa{و} \fa{پدرم} \fa{عصرها} \fa{تخته} \fa{نرد} \fa{بازی} \fa{میکنند}.
\end{quote}

\begin{quote}
\emph{pedarbozorg va pedaram 'asrhâ takhte nard bâzi mikonand.}
\end{quote}

\begin{quote}
Grandfather and my father play backgammon in the afternoons.
\end{quote}

\emph{takhte} 'board' + \emph{nard}. Often shortened to \emph{takhte}.

\begin{cousinweb}
\emph{takhte nard} joins inherited \emph{takhte}, which EtymDB traces to a Middle Persian word for 'tablet, plank, (chess) board' (EtymDB), and \emph{nard}, the game's own name, whose origin is not settled.
\end{cousinweb}

\begin{grammarlens}[title={What you've built}]
One more word for telling, explaining and describing at B1.
\end{grammarlens}

\subsection*{Your turn}
\begin{itemize}
\raggedright
  \item \emph{Say it:} \emph{takhte nard}
  \item \emph{Say it:} \emph{takhte nard}, once more
  \item \emph{Say it:} a sentence of your own with \emph{takhte nard}
  \item \emph{Recall:} say the Persian for pottery, then the Persian for a jigsaw puzzle, then say \emph{takhte nard} again
\end{itemize}

\begin{tcolorbox}[breakable,colback=teal!4,colframe=teal!35!black,title={Before you move on}]
What does \fa{تخته} \fa{نرد} mean? (\textquotedblleft{}Backgammon\textquotedblright{}.) Say it once more.
\end{tcolorbox}
\section[\texorpdfstring{\fa{ورق} (varaq)}{ورق (varaq)}]{\fa{ورق} (varaq) — a playing card; a sheet (of paper)}

\label{lesson:FA-C509-varaq}



\begin{quote}
Before the new one: say the Persian for a jigsaw puzzle, then the Persian for backgammon.
\end{quote}

\subsection*{You'll want to know: \fa{ورق}}
\textbf{\fa{ورق}} (\emph{varaq}) — \textquotedblleft{}a playing card; a sheet (of paper)\textquotedblright{}.

\begin{quote}
\fa{یک} \fa{دسته} \fa{ورق} \fa{در} \fa{کیف} \fa{من} \fa{بود}.
\end{quote}

\begin{quote}
\emph{yek daste varaq dar kif-e man bud.}
\end{quote}

\begin{quote}
There was a pack of cards in my bag.
\end{quote}

\emph{yek daste varaq} — a pack of cards. Also a sheet of paper or metal.

\begin{cousinweb}
\emph{varaq} comes from Arabic \emph{waraq} (EtymDB), a word for leaves and sheets of paper; Persian took the sheet and, from it, the playing card.
\end{cousinweb}

\begin{grammarlens}[title={What you've built}]
One more word for telling, explaining and describing at B1.
\end{grammarlens}

\subsection*{Your turn}
\begin{itemize}
\raggedright
  \item \emph{Say it:} \emph{varaq}
  \item \emph{Say it:} \emph{varaq}, once more
  \item \emph{Say it:} a sentence of your own with \emph{varaq}
  \item \emph{Recall:} say the Persian for a jigsaw puzzle, then the Persian for backgammon, then say \emph{varaq} again
\end{itemize}

\begin{tcolorbox}[breakable,colback=teal!4,colframe=teal!35!black,title={Before you move on}]
What does \fa{ورق} mean? (\textquotedblleft{}A playing card; a sheet (of paper)\textquotedblright{}.) Say it once more.
\end{tcolorbox}
\section[\texorpdfstring{\fa{تاس} (tâs)}{تاس (tâs)}]{\fa{تاس} (tâs) — a die (pl. dice)}

\label{lesson:FA-C509-tas}



\begin{quote}
Before the new one: say the Persian for backgammon, then the Persian for a playing card.
\end{quote}

\subsection*{You'll want to know: \fa{تاس}}
\textbf{\fa{تاس}} (\emph{tâs}) — \textquotedblleft{}a die (pl. dice)\textquotedblright{}.

\begin{quote}
\fa{او} \fa{تاس} \fa{را} \fa{انداخت} \fa{و} \fa{شش} \fa{آمد}.
\end{quote}

\begin{quote}
\emph{u tâs râ andâkht va shesh âmad.}
\end{quote}

\begin{quote}
He threw the die and got a six.
\end{quote}

\emph{tâs andâkhtan} or \emph{tâs rikhtan} — to roll the dice.

\begin{cousinweb}
\emph{tâs} is the Persian word for a die. EtymDB's only entry under this spelling is \emph{tâs}, 'cup', listed without a source, and the origin of the word for a die is not settled.
\end{cousinweb}

\begin{grammarlens}[title={What you've built}]
One more word for telling, explaining and describing at B1.
\end{grammarlens}

\subsection*{Your turn}
\begin{itemize}
\raggedright
  \item \emph{Say it:} \emph{tâs}
  \item \emph{Say it:} \emph{tâs}, once more
  \item \emph{Say it:} a sentence of your own with \emph{tâs}
  \item \emph{Recall:} say the Persian for backgammon, then the Persian for a playing card, then say \emph{tâs} again
\end{itemize}

\begin{tcolorbox}[breakable,colback=teal!4,colframe=teal!35!black,title={Before you move on}]
What does \fa{تاس} mean? (\textquotedblleft{}A die (pl. dice)\textquotedblright{}.) Say it once more.
\end{tcolorbox}
\section[\texorpdfstring{\fa{ورق} \fa{بازی} \fa{کردن} (varaq bâzi kardan)}{ورق بازی کردن (varaq bâzi kardan)}]{\fa{ورق} \fa{بازی} \fa{کردن} (varaq bâzi kardan) — to play cards}

\label{lesson:FA-C509-varaq-bazi-kardan}



\begin{quote}
Before the new one: say the Persian for a playing card, then the Persian for a die.
\end{quote}

\subsection*{You'll want to know: \fa{ورق} \fa{بازی} \fa{کردن}}
\textbf{\fa{ورق} \fa{بازی} \fa{کردن}} (\emph{varaq bâzi kardan}) — \textquotedblleft{}to play cards\textquotedblright{}.

\begin{quote}
\fa{شب} \fa{یلدا} \fa{با} \fa{فامیل} \fa{ورق} \fa{بازی} \fa{کردیم}.
\end{quote}

\begin{quote}
\emph{shab-e yaldâ bâ fâmil varaq bâzi kardim.}
\end{quote}

\begin{quote}
On Yalda night we played cards with the family.
\end{quote}

\emph{varaq} 'card' + \emph{bâzi kardan} 'to play'.

\begin{cousinweb}
\emph{varaq bâzi kardan} is built in Persian on \emph{varaq}, from Arabic \emph{waraq} (EtymDB), and the native phrase \emph{bâzi kardan}, 'to play'.
\end{cousinweb}

\begin{grammarlens}[title={What you've built}]
One more word for telling, explaining and describing at B1.
\end{grammarlens}

\subsection*{Your turn}
\begin{itemize}
\raggedright
  \item \emph{Say it:} \emph{varaq bâzi kardan}
  \item \emph{Say it:} \emph{varaq bâzi kardan}, once more
  \item \emph{Say it:} a sentence of your own with \emph{varaq bâzi kardan}
  \item \emph{Recall:} say the Persian for a playing card, then the Persian for a die, then say \emph{varaq bâzi kardan} again
\end{itemize}

\begin{tcolorbox}[breakable,colback=teal!4,colframe=teal!35!black,title={Before you move on}]
What does \fa{ورق} \fa{بازی} \fa{کردن} mean? (\textquotedblleft{}To play cards\textquotedblright{}.) Say it once more.
\end{tcolorbox}
